% Dirac/photon LSZ targets; detailed proofs are in the owning derivation file.
实验真正制备和探测的是进入相互作用区之前、离开相互作用区之后能够独立传播的粒子；而场论计算首先给出的却是若干时空点上的场关联函数。两者之间因此还缺少一步：必须从相关函数中取出与一个真实外部粒子对应的传播部分，只留下相互作用本身的振幅。

对稳定粒子，这一步可以从精确二点函数的质量壳结构看清楚。当四动量接近真实粒子的质量壳时，二点函数出现一个单粒子极点；极点的位置给出物理质量，极点的留数给出场算符与规范化单粒子态之间的重叠。LSZ 约化所做的就是用自由逆运动算符消去这段外部传播，再按极点留数归一化，并把结果投影到真实电子自旋态或真实光子的横向偏振态。得到的剩余部分才是散射振幅。

因此，“截肢外腿”不是在图上把最外面的一条线擦掉，而是对质量壳极点执行一个明确的代数留数运算。内部传播线并不代表渐近探测到的单粒子态，所以不进行同样的质量壳投影。


光子外腿使用物理四动量 $k^\mu=\hbar\kappa^\mu$。把\eqnrefs{eq:maxwell-kernel-projectors-r6} 的四波矢核改写成物理四动量核，定义
\begin{equation}
 \mathscr K_{\xi_{\rm gf}}^{\mu\nu}(k)
 \equiv
 \mathcal K_{\xi_{\rm gf}}^{\mu\nu}(\kappa=k/\hbar)
 =-\frac{k^2}{\mu_0\hbar^2}
 \left(P_T^{\mu\nu}+\frac1{\xi_{\rm gf}}P_L^{\mu\nu}\right).
 \label{eq:photon-physical-momentum-inverse-kernel-dp42}
\end{equation}
由\eqnrefs{eq:photon-prop-kappa-to-q-r13}，自由核与物理四动量表示的二点函数满足 $\mathscr K_{\mu\rho}\widetilde D_F^{\rho\nu}=\ii\hbar\,\delta_\mu{}^\nu$。电子传播子同样满足自由 Dirac 逆核乘传播子等于 $\ii\hbar$。因此两类外腿的截肢算符都必须含 $(\ii\hbar)^{-1}$：
\begin{equation}
 \boxed{
 \begin{aligned}
 \text{电子外腿:}&\quad
 \widetilde S(p)\ \xrightarrow{\;[Z_{\rm pole}^{(e)}]^{-1/2}(\slashed p-mc)/(\ii\hbar)\;}\
 u,\bar u,v,\bar v,\\
 \text{光子外腿:}&\quad
 \widetilde D_{\mu\nu}(k)\
 \xrightarrow{\;[Z_{\rm pole}^{(\gamma)}]^{-1/2}\mathscr K^{\mu\rho}(k)/(\ii\hbar)\;}\
 \varepsilon_\rho^{(\lambda)}(k),
 \qquad k\cdot\varepsilon^{(\lambda)}=0.
 \end{aligned}}
 \label{eq:lsz-dirac-photon-mainline-dp8}
\end{equation}
这里用 $Z_{\rm pole}^{(e)}$ 与 $Z_{\rm pole}^{(\gamma)}$ 专门表示当前所用场变量的电子与横向光子单粒子极点权重。它们是二点函数的谱数据，负责把局域场算符产生的状态归一化成单位范数渐近粒子。重整化中常用的 $Z_2,Z_3$ 则表示“裸场与重整化场之间的场强重整化常数”；在一般方案中，它们与单粒子极点权重是不同对象。若采用在壳方案，$Z_2,Z_3$ 被选择来使重整化传播子的相应物理极点留数等于 $1$。

\begin{figure}[htbp]
 \centering
 \begin{tikzpicture}[x=1cm,y=1cm,every node/.style={font=\large},>=Stealth]
 \node[font=\bfseries] at (-3.5,1.55) {(a) full Green function};
 \node[font=\bfseries] at (3.55,1.55) {(b) amputated pole};

 % electron row
 \draw[fermion] (-5.3,.55)--(-2.7,.55);
 \node[above] at (-4.0,.55) {$S(p)$};
 \node[draw,rounded corners=3pt,minimum width=1.45cm,minimum height=.9cm] at (-1.95,.55) {$\Gamma$};
 \node at (0,.55) {$\xrightarrow{\quad Z_2^{-1/2}(\slashed p-mc)\quad}$};
 \node at (2.15,.55) {$\bar u_s(p)$};
 \draw[fermion] (2.75,.55)--(3.35,.55);
 \node[draw,rounded corners=3pt,minimum width=1.45cm,minimum height=.9cm] at (4.1,.55) {$\Gamma$};

 % photon row
 \draw[photon] (-5.3,-1.0)--(-2.7,-1.0);
 \node[above] at (-4.0,-1.0) {$D_{\mu\nu}(k)$};
 \node[draw,rounded corners=3pt,minimum width=1.45cm,minimum height=.9cm] at (-1.95,-1.0) {$\Gamma^\nu$};
 \node at (0,-1.0) {$\xrightarrow{\quad Z_3^{-1/2}\mathcal K^{\rho\mu}_\xi\quad}$};
 \node at (2.05,-1.0) {$\varepsilon_\rho^{(\lambda)*}$};
 \draw[photon] (2.85,-1.0)--(3.35,-1.0);
 \node[draw,rounded corners=3pt,minimum width=1.45cm,minimum height=.9cm] at (4.1,-1.0) {$\Gamma^\rho$};
\end{tikzpicture}

 \caption{LSZ 约化的图形表示。完整 Green 函数含外腿单粒子极点；自由逆运动核除以 $\ii\hbar$ 后消去该传播极点，再投影到壳上电子自旋子或横向光子偏振。$\Gamma$ 表示其余连通截肢核。}
 \label{fig:qed-lsz-amputation-dp8}
\end{figure}

电子场与物理单电子态的重叠定义电子极点留数：
\begin{equation}
 \boxed{
 \langle\Omega|\psi(0)|p,s\rangle
 =\sqrt{Z_{\rm pole}^{(e)}}\,u_s(p),}
 \qquad
 (\slashed p-m_{\rm phys}c)u_s(p)=0.
 \label{eq:dirac-field-pole-overlap-dp8}
\end{equation}
于是精确电子二点函数在质量壳附近具有
\begin{equation}
 \boxed{
 \widetilde S(p)
 \xrightarrow[p^2\to m_{\rm phys}^2c^2]{}
 \frac{\ii\hbar Z_{\rm pole}^{(e)}(\slashed p+m_{\rm phys}c)}
 {p^2-m_{\rm phys}^2c^2+\ii0^+}
 +\widetilde S_{\rm reg}(p).}
 \label{eq:dirac-exact-pole-lsz-dp8}
\end{equation}
$\widetilde S_{\rm reg}$ 在该单粒子极点邻域解析，并包含多粒子连续谱贡献。若连通 Green 函数的一条出射电子外腿在该极点附近写成 $\widetilde G(p;\{q\})=\widetilde S(p)\widetilde\Gamma(p;\{q\})+\text{非极点项}$，则电子外腿的留数操作为
\begin{equation}
 \boxed{
 \lim_{p^2\to m_{\rm phys}^2c^2}
 \bar u_s(p)[Z_{\rm pole}^{(e)}]^{-1/2}
 \frac{\slashed p-m_{\rm phys}c}{\ii\hbar}
 \widetilde G(p;\{q\})
 =\bar u_s(p)\widetilde\Gamma(p;\{q\}).}
 \label{eq:dirac-lsz-amputation-r68}
\end{equation}
这条式子把“二点函数具有单粒子极点”转换成“外腿可以被截肢并投影到壳上自旋子”。电子与正电子的四种渐近外态只是同一极点投影在交叉变换下的不同取向：
\begin{equation}
 \begin{array}{c|c}
 \text{渐近外态} & \text{壳上波函数}\\
 \hline
 \text{入射 }e^- & u_s(p)\\
 \text{出射 }e^- & \bar u_s(p)\\
 \text{入射 }e^+ & \bar v_s(p)\\
 \text{出射 }e^+ & v_s(p)
 \end{array}
 \label{eq:external-dirac-spinor-table-dp8}
\end{equation}

光子外态满足同样的极点逻辑，但只在两个横向物理偏振上取留数。定义物理光子极点权重
\begin{equation}
 \langle\Omega|A_\mu(0)|k,\lambda\rangle
 =\sqrt{Z_{\rm pole}^{(\gamma)}}\,\varepsilon_\mu^{(\lambda)}(k),
 \qquad
 k^2=0,
 \qquad
 k\cdot\varepsilon^{(\lambda)}=0.
 \label{eq:photon-field-pole-overlap-r68}
\end{equation}
若一条光子外腿的连通 Green 函数写成 $\widetilde G_\nu(k;\{q\})=\widetilde D_{\nu\rho}(k)\widetilde\Gamma^\rho(k;\{q\})+\text{非极点项}$，则相应的留数操作为
\begin{equation}
 \boxed{
 \lim_{k^2\to0}
 \varepsilon_\mu^{(\lambda)*}(k)[Z_{\rm pole}^{(\gamma)}]^{-1/2}
 \frac{\mathscr K_{\xi_{\rm gf}}^{\mu\nu}(k)}{\ii\hbar}
 \widetilde G_\nu(k;\{q\})
 =\varepsilon_\rho^{(\lambda)*}(k)\widetilde\Gamma^\rho(k;\{q\}).}
 \label{eq:photon-lsz-amputation-r68}
\end{equation}
规范固定核的纵向部分可以出现在中间传播子中；与物理偏振或守恒流收缩后，它不进入散射矩阵元。对每一条外腿重复同一留数操作，连通 Green 函数的最高阶多重极点便分解为外腿极点与一个共同的截肢核，最终得到
\begin{equation}
 \boxed{
 \mathcal M_{fi}
 =\left[\prod_{\rm ext\;fermions}u,\bar u,v,\bar v\right]
 \left[\prod_{\rm ext\;photons}\varepsilon^{(\lambda)}\right]
 \widetilde\Gamma_{\rm amp}.}
 \label{eq:multi-leg-lsz-r68}
\end{equation}
因此电子与光子的 LSZ 约化共同实现“精确单粒子极点 $\rightarrow$ 渐近物理态 $\rightarrow$ 连通截肢振幅”的映射。
